\section{Categorical logic}
\label{sec:type-theory}

We now give an internal language for representable
$\catname{SLat}$-multicategories.  Following Shulman~\cite{shulman2016categorical},
we let inference rules absorb categorical structure: contexts are the inputs
of a multimorphism and substitution is composition.  Detailed coherence and
substitution arguments appear in~\Cref{app:proofs}.

\begin{definition}
A \emph{multigraph} $\mathcal{G}$ consists of a set of objects and, for every list $A_1,\ldots,A_n$ and object $B$, a set $\mathcal{G}(A_1,\ldots,A_n;B)$ of generating arrows.  We regard a multigraph as a many-sorted signature.
\end{definition}

\begin{definition}[The core calculus]
\label{def:s_lat_type_theory}
A context is a finite ordered list of pairwise distinct typed variables;
$()$ denotes the empty context, and commas denote concatenation.  The type
theory generated by a multigraph $\mathcal{G}$ has the following rules:
\begin{gather*}
  \inferrule*{A\in\obj{\mathcal G}}{\vdash A\;\mathrm{type}}
  \qquad
  \inferrule*{\vdash A\;\mathrm{type}\quad\vdash B\;\mathrm{type}}{\vdash A\otimes B\;\mathrm{type}}
  \qquad
  \inferrule*{\;}{\vdash\I\;\mathrm{type}}
  \qquad
  \inferrule*[right=$\id$]{\vdash A\;\mathrm{type}}{x:A\vdash x:A}
  \\
  \inferrule*[right=$f$]
  {f\in\mathcal G(A_1,\ldots,A_n;B)
   \quad \Gamma_l\vdash\sum_{j=1}^{i_l}M_{lj}:A_l\ (1\leq l\leq n)}
  {\Gamma_1,\ldots,\Gamma_n\vdash
   \sum_{k_1=1}^{i_1}\cdots\sum_{k_n=1}^{i_n}
   f(M_{1k_1},\ldots,M_{nk_n}):B}
  \\
  \inferrule*[right=$\otimes_{\mathrm{intro}}$]
  {\Gamma\vdash\sum_{i=1}^{m}M_i:A\quad
   \Delta\vdash\sum_{j=1}^{n}N_j:B}
  {\Gamma,\Delta\vdash\sum_{i=1}^{m}\sum_{j=1}^{n}\{M_i,N_j\}:A\otimes B}
  \\
  \inferrule*[right=$\otimes_{\mathrm{elim}}$]
  {\Gamma\vdash\sum_{i=1}^{m}M_i:A\otimes B\quad
   \Delta_1,x:A,y:B,\Delta_2\vdash\sum_{j=1}^{n}N_j:C}
  {\Delta_1,\Gamma,\Delta_2\vdash
   \sum_{i=1}^{m}\sum_{j=1}^{n}\text{\rm match}(M_i,xy.N_j):C}
  \\
  \inferrule*[right=$\I_{\mathrm{intro}}$]{\;}{()\vdash\star:\I}
  \qquad
  \inferrule*[right=$\I_{\mathrm{elim}}$]
  {\Gamma\vdash\sum_{i=1}^{m}M_i:\I\quad
   \Delta_1,\Delta_2\vdash\sum_{j=1}^{n}N_j:A}
  {\Delta_1,\Gamma,\Delta_2\vdash
   \sum_{i=1}^{m}\sum_{j=1}^{n}\text{\rm match}_{\I}(M_i,N_j):A}
  \\
  \inferrule*[right=$\join_{\mathrm{intro}}$]
  {\Gamma\vdash M:A\quad\Gamma\vdash N:A}
  {\Gamma\vdash M\join N:A}.
\end{gather*}
Here and below,
\(
  \sum_{i=1}^{n}M_i
\)
is notation for the iterated join
\(M_1\join \cdots\join M_n\), under some bracketing; it is not an additional term constructor.  
The displayed summands in the premises are join-free.  Thus every conclusion is a finite non-empty join of join-free terms, with joins occurring only at the top level.  Derivations, rather than independently generated raw terms, are the primary syntax.
When a constructor is displayed with iterated joins in its arguments, we use
the corresponding \emph{extended application}: the notation abbreviates the
distributed conclusion of its typing rule.
In $\text{\rm match}(M,xy.N)$, the variables $x$ and $y$ are bound in $N$;
all substitutions below are capture-avoiding.
\end{definition}

\begin{definition}[Notation for derivations and joins]
If $\mathcal D$ is a derivation of $\Gamma\vdash M:A$, we write
$\mathcal D::\Gamma\vdash M:A$ and $|\mathcal D|=M$.  A derivable term is
\emph{join-free} if it contains no occurrence of $\join$.  Flattening the
top-level joins of a derivable term $M$ gives a finite non-empty list of
join-free terms, called its \emph{summands}.  We write
$\operatorname{supp}(M)$ for the set of these summands, forgetting their order
and multiplicity.  For derivable terms with the same context and result type,
we define
\[
M\equiv_{\join}N
\quad\Longleftrightarrow\quad
\operatorname{supp}(M)=\operatorname{supp}(N).
\]
Equivalently, $\equiv_{\join}$ is the equivalence generated by associativity,
commutativity, and idempotence of $\join$.

For derivations with the same context and result type (but not necessarily the
same raw conclusion), $\mathcal D\sim\mathcal E$ denotes
\emph{permutation equivalence}: the equivalence generated by the corresponding
semilattice rearrangements of $\join_{\mathrm{intro}}$ and by commuting
$\join_{\mathrm{intro}}$ past the other inference rules.  Its precise
generators are given in~\Cref{def:perm_equiv}.
\end{definition}

The intended construction has types as the objects of a free representable
multicategory enriched in $\catname{SLat}$, and derivations as its morphisms.  We must
therefore equip derivations with a composition operation.  This operation is
substitution.  To define multicategorical composition, substitution must be
unital, associative, and satisfy interchange
(\Cref{remark:placed_composition}); to define composition in an
$\catname{SLat}$-multicategory, it must in addition preserve finite non-empty
joins in each argument.  These requirements motivate both the rules above and
the results below.

The distributive conclusions make an application to joins produce the join of
all choices of representatives.  Thus syntax absorbs multilinearity instead
of imposing it separately for every constructor.  The price is that a
conclusion term need not determine its derivation.  The
coherence theorem in~\Cref{app:core-coherence} says precisely that
\[
|\mathcal D|\equiv_{\join}|\mathcal E|
\quad\Longleftrightarrow\quad
\mathcal D\sim\mathcal E
\qquad\text{(\Cref{theorem:coherence}).}
\]
Given derivations $\mathcal D::\Gamma\vdash M:A$ and
$\mathcal E::\Delta_1,x:A,\Delta_2\vdash N:B$, substitution is defined by
induction on $\mathcal E$ as a derivation
\[
\mathsf{sub}_x(\mathcal E,\mathcal D)
::\Delta_1,\Gamma,\Delta_2\vdash N[M/x]:B.
\]
The notation $N[M/x]$ denotes its conclusion up to $\equiv_{\join}$; the
independence of the chosen derivations is part of the result below.
Intuitively, it substitutes every summand of $M$ into every summand of $N$:
\[
\operatorname{supp}(N[M/x])
=
\{\,q[m/x]\mid q\in\operatorname{supp}(N),\
m\in\operatorname{supp}(M)\,\}.
\]
\begin{lemma}[Admissibility of substitution]
\label{lemma:subst_admissible_main}
If $\Gamma\vdash M:A$ and
$\Delta_1,x:A,\Delta_2\vdash N:B$ are derivable, then so is
\[
\Delta_1,\Gamma,\Delta_2\vdash N[M/x]:B.
\]
\end{lemma}

\begin{lemma}[Properties of substitution]
\label{lemma:subst_properties_main}
Substitution is unital and preserves joins in both arguments.  Moreover, for
well-typed terms it satisfies associativity and interchange:
\begin{align*}
(P[M/x])[N/y]&\equiv_{\join}P[M[N/y]/x],\\
P[M/x][N/y]&\equiv_{\join}P[N/y][M/x],\\
\left(\sum_{i=1}^{n}N_i\right)
  \left[\sum_{j=1}^{m}M_j/x\right]
&\equiv_{\join}
  \sum_{i=1}^{n}\sum_{j=1}^{m}N_i[M_j/x].
\end{align*}
\end{lemma}

Coherence makes the conclusion of derivation-level substitution independent
of its chosen derivations; see \Cref{def:substitution},
\Cref{lemma:subst_admissible,lemma:subst_properties}, and
\Cref{prop:subst_derivation_independence}.
Hence non-unique derivations induce the same composite after quotienting by
the semilattice laws.

\begin{definition}[Equality]
\label{def:slat_equality}
Equality $\Gamma\vdash M\equiv N:A$ is the least congruence on derivable terms containing
\begin{align*}
M\join(N\join P)&\equiv(M\join N)\join P,
&M\join N&\equiv N\join M,\\
M\join M&\equiv M,\\
\text{match}(\{P,Q\},xy.R)&\equiv R[P/x][Q/y],\\
\text{match}(S,xy.R[\{x,y\}/u])&\equiv R[S/u],\\
\text{match}_{\I}(\star,R)&\equiv R,
&\text{match}_{\I}(S,R[\star/u])&\equiv R[S/u].
\end{align*}
The first three equations give the associative, commutative, and idempotent
laws for $\join$; the remaining equations express the universal properties of
$\otimes$ and $\I$.  Equality is closed under all term constructors, and
stability under substitution is admissible.  The complete rule schemes and
the proof that derivation-level substitution descends to $\equiv$-classes
appear in \Cref{def:slat_equality_details} and
\Cref{prop:subst_respects_equiv}.
\end{definition}

\begin{theorem}[Initiality]
\label{theorem:free_slat_multicat}
For every multigraph $\mathcal G$, the types and equivalence classes of derivable judgments of~\Cref{def:s_lat_type_theory} present the free representable $\catname{SLat}$-multicategory on $\mathcal G$.  Equivalently, they compute the left adjoint in
\[
\catname{Multigraph}
\;\underset{U}{\stackrel{F}{\rightleftarrows}}\;
\catname{SLat}\text{-}\catname{RepMultCat}.
\]
\end{theorem}
Here the representable multicategories carry chosen representations, and the
morphisms preserve those choices.  The proof is deferred
to~\Cref{app:initiality-details}.

Ordinary e-graphs use the cartesian $\mathfrak S$-multicategory in which
$\mathfrak S$ contains all functions between finite cardinals.  Following the
structural type theories of Shulman~\cite{shulman2016categorical}, we replace
the linear rules by their common-context versions: every premise of a
term-forming rule has the same ambient context.  In particular, the identity
and unit rules become
\begin{equation}
\label{eq:cartesian_structure}
\inferrule*[right=$\id$]{x:A\in\Gamma}{\Gamma\vdash x:A}
\qquad
\inferrule*{\;}{\Gamma\vdash\star:\I}.
\end{equation}
For example, every premise of the $f$-rule now has context $\Gamma$:
\[
\inferrule*[right=$f$]
 {f\in\mathcal G(A_1,\ldots,A_n;B)\quad
  \Gamma\vdash\sum_j M_{lj}:A_l\ (1\leq l\leq n)}
 {\Gamma\vdash
  \sum_{k_1,\ldots,k_n}f(M_{1k_1},\ldots,M_{nk_n}):B}.
\]
For a nullary $f$, the ambient context $\Gamma$ is arbitrary.  The other rules
are changed in the same way.  Consequently variables may be
exchanged, copied, or discarded.  Substitution is correspondingly
simultaneous rather than linear: from
\(\Gamma\vdash M_i:A_i\) and
\(\Gamma,\vec x:\vec A\vdash N:B\) one obtains
\(\Gamma\vdash N[\vec M/\vec x]:B\).
Here $\vec x:\vec A$ abbreviates
$x_1:A_1,\ldots,x_n:A_n$, and $[\vec M/\vec x]$ denotes simultaneous
capture-avoiding substitution.

\begin{lemma}[Cartesian substitution]
\label{lemma:cartesian_substitution_main}
Simultaneous substitution is admissible and preserves finite non-empty joins.
Its associativity and interchange laws combine into the usual cartesian
substitution law
\[
P[N/y][M/x]\equiv_{\join}P[M/x][N[M/x]/y],
\]
when $\Gamma\vdash M:A$, $\Gamma,x:A\vdash N:B$, and
$\Gamma,x:A,y:B\vdash P:C$ (up to fresh renaming of bound variables).
Substitution also commutes with the $\mathfrak S$-action, hence with
weakening, exchange, and contraction.
\end{lemma}
These are the standard structural substitution laws; see
Shulman~\cite{shulman2016categorical} and \Cref{app:multicategories}.

\begin{definition}[Presentations]
\label{def:type_theory_under_presentation}
A multicategorical presentation is a pair $\mathcal P=(\mathcal G,\mathcal E)$ of a multigraph and a set of parallel generating equations.  We write
$\mathcal E(\Gamma;B)$ for the equations whose two sides have context $\Gamma$
and type $B$.  The type theory under $\mathcal P$
extends~\Cref{def:s_lat_type_theory} by
\[
\inferrule*{(M,N)\in\mathcal E(\Gamma;B)}
{\Gamma\vdash M\equiv N:B}.
\]
Its quotient presents the free representable $\catname{SLat}$-multicategory satisfying $\mathcal E$.
\end{definition}

For the running e-graph, $\mathcal G$ contains the operations $*$, $/$, and $<\!\!<$ and the constants, while $\mathcal E$ contains, among others,
\[
x*2\equiv x<\!\!<1,
\qquad
(x*y)/z\equiv x*(y/z),
\qquad
x/x\equiv1,
\qquad
x*1\equiv x.
\]
Thus equality saturation can be read as rewriting modulo the structural equality above, while $\join$ records the representatives accumulated in an e-class.

\subsection{Explicit sharing and recursion}

\begin{definition}[Explicit sharing]
\label{def:explicit_sharing}
The admissible cut can be exposed as a term former
\[
\inferrule*[right=let]
{\Gamma\vdash M:A\quad\Gamma,x:A\vdash N:B}
{\Gamma\vdash
 \textbf{let }x=M\textbf{ in }N:B},
\qquad
\textbf{let }x=M\textbf{ in }N\equiv N[M/x].
\]
Congruence is imposed in both positions.
\end{definition}
This is an extension of the cartesian $\mathfrak S$-calculus above.  In
particular, the two premises have the same ambient context, and eliminating a
let uses the simultaneous substitution of
\Cref{lemma:cartesian_substitution_main}.  The extension is definitional:
eliminating all lets returns a core term and does not change the presented
category.  Nevertheless, retaining a let records sharing.  In particular, a
join in an argument position is represented by binding the whole e-class
rather than duplicating it.

\begin{example}[The acyclic e-graph]
\label{ex:acyclic_let}
Example~(b) of~\Cref{fig:e-graph-example} is represented by
\[
\vdash
\textbf{let }x=a\textbf{ in }
\textbf{let }z=2\textbf{ in }
\textbf{let }e=x*z\join x\mathbin{<\!\!<}1\textbf{ in }e/z
:\mathbb N.
\]
Here $e$ names the e-class $\{x*z,x\mathbin{<\!\!<}1\}$, while $z$ is shared by the multiplication and division nodes.  Expanding the lets gives
$(a*2)/2\join(a\mathbin{<\!\!<}1)/2$.
For $c\in\{a,1,2\}$ and any context $\Gamma$, write $\kappa_c^\Gamma$
for the nullary instance of the $f$-rule
\[
\inferrule*[left=$\kappa_c^\Gamma$]
 {c\in\mathcal G(;\mathbb N)}
 {\Gamma\vdash c:\mathbb N}.
\]
Together with the cartesian identity rule, these close all the leaves.  The
complete derivation follows; every triangle refers to a subderivation displayed
above it.
{\footnotesize
\setlength{\jot}{0pt}
\begin{gather*}
\inferrule*[left=$\pi_{x*z}$]
 {x:\mathbb N,z:\mathbb N\vdash x:\mathbb N\quad
  x:\mathbb N,z:\mathbb N\vdash z:\mathbb N\quad
  {*}\in\mathcal G(\mathbb N,\mathbb N;\mathbb N)}
 {x:\mathbb N,z:\mathbb N\vdash x*z:\mathbb N}
\\
\inferrule*[left=$\pi_{x\mathbin{<\!\!<}1}$]
 {x:\mathbb N,z:\mathbb N\vdash x:\mathbb N\quad
  \pitri{\kappa_1^{x,z}}\quad
  {<\!\!<}\in\mathcal G(\mathbb N,\mathbb N;\mathbb N)}
 {x:\mathbb N,z:\mathbb N\vdash x\mathbin{<\!\!<}1:\mathbb N},\\
\inferrule*[right=$\join_{\mathrm{intro}}$,left=$\pi_E$]
 {\pitri{\pi_{x*z}}\quad\pitri{\pi_{x\mathbin{<\!\!<}1}}}
 {x:\mathbb N,z:\mathbb N\vdash x*z\join x\mathbin{<\!\!<}1:\mathbb N}
\\
\inferrule*[left=$\pi_{e/z}$]
 {x:\mathbb N,z:\mathbb N,e:\mathbb N\vdash e:\mathbb N\quad
  x:\mathbb N,z:\mathbb N,e:\mathbb N\vdash z:\mathbb N\quad
  {/}\in\mathcal G(\mathbb N,\mathbb N;\mathbb N)}
 {x:\mathbb N,z:\mathbb N,e:\mathbb N\vdash e/z:\mathbb N},\\
\inferrule*[right=let,left=$\pi_e$]
 {\pitri{\pi_E}\quad\pitri{\pi_{e/z}}}
 {x:\mathbb N,z:\mathbb N\vdash
  \textbf{let }e=x*z\join x\mathbin{<\!\!<}1\textbf{ in }e/z:\mathbb N},\\
\inferrule*[right=let,left=$\pi_z$]
 {\pitri{\kappa_2^x}\quad\pitri{\pi_e}}
 {x:\mathbb N\vdash
  \textbf{let }z=2\textbf{ in }
  \textbf{let }e=x*z\join x\mathbin{<\!\!<}1\textbf{ in }e/z:\mathbb N}.
\end{gather*}
\[
\adjustbox{max width=\textwidth}{$
\inferrule*[right=let]
 {\pitri{\kappa_a^\varnothing}\quad\pitri{\pi_z}}
 {\vdash
  \textbf{let }x=a\textbf{ in }
  \textbf{let }z=2\textbf{ in }
  \textbf{let }e=x*z\join x\mathbin{<\!\!<}1\textbf{ in }e/z:\mathbb N}
$}.
\]
}
\end{example}

Acyclic sharing is not sufficient for example~(e), whose root e-class is used recursively.  We therefore add mutually recursive bindings.

\begin{definition}[Recursive sharing]
\label{def:recursive_sharing}
A declaration list is $D=(x_1=M_1,\ldots,x_n=M_n)$ with distinct $x_i$.  Recursive sharing is typed by
\[
\inferrule*[right=letrec]
{\Gamma,\vec x:\vec A\vdash M_i:A_i\ (1\leq i\leq n)
 \quad
 \Gamma,\vec x:\vec A\vdash\sum_{k=1}^{m}N_k:B}
{\Gamma\vdash
 \sum_{k=1}^{m}\textbf{letrec }D\textbf{ in }N_k:B}.
\]
The declarations and body have the same ambient context, so the $M_i$ may be mutually recursive.  Joins distribute over the body, but not over binding right-hand sides.  The calculus is quotiented by unfolding, merging, scope extrusion, garbage collection, unpairing, and diagonal equations; their precise formulation is~\Cref{eq:let_equations}.
\end{definition}
The design translates the symmetric monoidal trace axioms into multicategorical
form and draws on Hasegawa's models of sharing graphs~\cite{hasegawa_models_1999};
e-graphs are reminiscent of sharing graphs with the addition of e-classes.

For reference, a multicategorical trace on a representable multicategory
$\mathcal M$ is a family of operations
\[
\operatorname{Tr}_{\Gamma,B}^{X}:
\mathcal M(\Gamma,X;B\otimes X)\longrightarrow\mathcal M(\Gamma;B)
\]
that feeds the $X$-component of the output back into the $X$-input and satisfies
the usual tightening, sliding, vanishing, strength, and yanking axioms.  The
full axioms are recalled in~\Cref{app:sharing-trace-details}.
The restriction on recursive right-hand sides is essential.  Enrichment makes composition multilinear, but does not make trace additive.  For example, in $\catname{Rel}$ let $A=\{a\}$, $B=\{b\}$, and $X=\{0,1\}$, and consider relations from $A\sqcup X$ to $B\sqcup X$:
\[
R_1=\{(a,0),(1,b)\},
\qquad
R_2=\{(0,1)\}.
\]
Neither relation separately has a feedback path from $a$ to $b$, so both traces are empty, whereas
$R_1\cup R_2$ contains $a\to0\to1\to b$.  Hence $(a,b)\in\operatorname{Tr}(R_1\cup R_2)$ but
$(a,b)\notin\operatorname{Tr}(R_1)\cup\operatorname{Tr}(R_2)$.
A join inside a recursive declaration must therefore remain protected.  For an actual e-class whose members are already provably equal, the expected distribution is derivable using congruence and idempotence.

This protection also drives coherence.  Away from a binder,
$\join_{\mathrm{intro}}$ commutes towards the root as before; in a binding
right-hand side, a join remains part of one atomic binding.  Derivations thus
decompose into finite joins of binding-atomic components.  Substitution
distributes over exposed components but treats an atomic binding as a whole,
avoiding the mixed feedback loops caused by full distribution.

Write $\llbracket\Gamma\vdash M:B\rrbracket$ for the morphism represented by
the equivalence class of the displayed judgment.  The recursive binder induces
the multicategorical trace
\begin{equation}
\operatorname{Tr}_{\Gamma,B}^{X}
\bigl(\llbracket\Gamma,x:X\vdash M:B\otimes X\rrbracket\bigr)
=
\llbracket
\Gamma\vdash
\textbf{letrec }x=\text{\rm match}(M,uv.v)
\textbf{ in }\text{\rm match}(M,uv.u):B
\rrbracket.
\label{eq:syntactic_trace}
\end{equation}

\begin{proposition}
\label{prop:traced_representable}
There is a $2$-equivalence between traced symmetric representable
multicategories and traced symmetric monoidal categories.
\end{proposition}

\begin{proposition}[Soundness of the syntactic trace]
\label{prop:trace_construction}
The operation in~\eqref{eq:syntactic_trace}, together with the equations of~\Cref{eq:let_equations}, satisfies the multicategorical trace axioms.
\end{proposition}

Proofs of the preceding two propositions appear in~\Cref{app:derivations}.

\begin{theorem}[Initiality of recursive sharing]
\label{theorem:free_traced_slat}
For every multigraph $\mathcal G$, the cartesian calculus generated by $\mathcal G$ and extended with \textbf{letrec} presents the free traced cartesian representable $\catname{SLat}$-multicategory on $\mathcal G$.  Equivalently, it presents the free traced cartesian monoidal $\catname{SLat}$-category on $\mathcal G$.
\end{theorem}
The protected-join coherence and substitution results, followed by the
initiality proof, appear in
\Cref{def:binding_atomic,def:letrec_substitution},
\Cref{theorem:binding_coherence,theorem:letrec_subst_properties,theorem:letrec_freeness_details},
and \Cref{prop:letrec_subst_respects_equiv}.

\begin{example}[The cyclic e-graph]
\label{ex:cyclic_letrec}
Example~(e) of~\Cref{fig:e-graph-example} is represented extensionally by
\[
\vdash
\textbf{letrec }x=
 (x*2\join x\mathbin{<\!\!<}1)/2\join x*(2/2\join 1)\join a
\textbf{ in }x:\mathbb N.
\]
A complete derivation follows; every triangle refers to a subderivation
displayed above it.
{\footnotesize
\setlength{\jot}{0pt}
\begin{gather*}
\inferrule*[left=$\pi_{x*2}$]
 {x:\mathbb N\vdash x:\mathbb N\quad
  \pitri{\kappa_2^x}\quad
  {*}\in\mathcal G(\mathbb N,\mathbb N;\mathbb N)}
 {x:\mathbb N\vdash x*2:\mathbb N}
\\
\inferrule*[left=$\pi_{x\mathbin{<\!\!<}1}$]
 {x:\mathbb N\vdash x:\mathbb N\quad
  \pitri{\kappa_1^x}\quad
  {<\!\!<}\in\mathcal G(\mathbb N,\mathbb N;\mathbb N)}
 {x:\mathbb N\vdash x\mathbin{<\!\!<}1:\mathbb N},\\
\inferrule*[right=$\join_{\mathrm{intro}}$,left=$\pi_E$]
 {\pitri{\pi_{x*2}}\quad\pitri{\pi_{x\mathbin{<\!\!<}1}}}
 {x:\mathbb N\vdash x*2\join x\mathbin{<\!\!<}1:\mathbb N}
\\
\inferrule*[left=$\pi_{E/2}$]
 {\pitri{\pi_E}\quad\pitri{\kappa_2^x}\quad
  {/}\in\mathcal G(\mathbb N,\mathbb N;\mathbb N)}
 {x:\mathbb N\vdash
  (x*2\join x\mathbin{<\!\!<}1)/2:\mathbb N},\\
\inferrule*[right=$\join_{\mathrm{intro}}$,left=$\pi_W$]
 {\inferrule*{\pitri{\kappa_2^x}\quad
               \pitri{\kappa_2^x}\quad
               {/}\in\mathcal G(\mathbb N,\mathbb N;\mathbb N)}
              {x:\mathbb N\vdash2/2:\mathbb N}
  \quad\pitri{\kappa_1^x}}
 {x:\mathbb N\vdash2/2\join 1:\mathbb N}
\\
\inferrule*[left=$\pi_{x*W}$]
 {x:\mathbb N\vdash x:\mathbb N\quad\pitri{\pi_W}\quad
  {*}\in\mathcal G(\mathbb N,\mathbb N;\mathbb N)}
 {x:\mathbb N\vdash x*(2/2\join 1):\mathbb N},\\
\inferrule*[right=$\join_{\mathrm{intro}}$,left=$\pi_Z'$]
 {\pitri{\pi_{E/2}}\quad\pitri{\pi_{x*W}}}
 {x:\mathbb N\vdash
  (x*2\join x\mathbin{<\!\!<}1)/2
  \join x*(2/2\join 1):\mathbb N},\\
\inferrule*[right=$\join_{\mathrm{intro}}$,left=$\pi_Z$]
 {\pitri{\pi_Z'}\quad\pitri{\kappa_a^x}}
 {x:\mathbb N\vdash
  (x*2\join x\mathbin{<\!\!<}1)/2
  \join x*(2/2\join 1)\join a:\mathbb N}.
\end{gather*}
\[
\adjustbox{max width=\textwidth}{$
\inferrule*[right=letrec]
 {\pitri{\pi_Z}\quad x:\mathbb N\vdash x:\mathbb N}
 {\vdash
  \textbf{letrec }x=
  (x*2\join x\mathbin{<\!\!<}1)/2
  \join x*(2/2\join 1)\join a
  \textbf{ in }x:\mathbb N}
$}.
\]
}
Its sharing structure can be shown explicitly as
\[
\begin{aligned}
\vdash\;&\textbf{let }z=2\textbf{ in }
\textbf{let }w=1\textbf{ in }\\
&\textbf{letrec }x=
 \bigl(\textbf{let }e=x*z\join x\mathbin{<\!\!<}w\textbf{ in }e/z\bigr)\\
&\hspace{5.7em}\join
 \bigl(\textbf{let }c=z/z\join w\textbf{ in }x*c\bigr)\join a
\textbf{ in }x:\mathbb N.
\end{aligned}
\]
The recursive variable $x$ names the cyclic root class, while $e$ and $c$ name the two shared inner e-classes.
\end{example}

\begin{remark}
\label{remark:no_free_enrichment_factorisation}
Unlike \textbf{let}, \textbf{letrec} is not eliminated by a terminating definitional equation: unfolding may continue indefinitely.  
It therefore adds genuine traced morphisms rather than merely a sharing notation for morphisms of the core calculus.  
This is also why the recursive calculus must be constructed directly rather than obtained by freely enriching an ordinary traced category.
\end{remark}
